\documentclass[a4paper,12pt]{article}
\usepackage[utf8]{inputenc}
\usepackage{hyperref}
\usepackage{array}
\usepackage{lastpage}
\usepackage{lipsum}
\usepackage{boldline} 
\usepackage{tabularx, colortbl}
\usepackage{fancyhdr}
\usepackage{xltabular}
\usepackage{amssymb}
\usepackage{graphicx}
\usepackage[hmargin=2cm,top=5.5cm,headheight=130pt,footskip=65pt]{geometry}
\usepackage{amsthm}

\theoremstyle{definition}
\newtheorem{theorem}{REQ}


\newcount\savedtheorem
%\stepcounter{theorem}
\newenvironment{subtheorems}{%
  \savedtheorem=\value{theorem}%
  \edef\prevthetheorem{\thetheorem}%
  \setcounter{theorem}{0}%
  \renewcommand\thetheorem{\prevthetheorem.\arabic{theorem}}%
}
{%
  \setcounter{theorem}{\savedtheorem}%
}

\renewcommand{\familydefault}{\sfdefault}

\hypersetup{
    colorlinks=true,
    linkcolor=blue,
    filecolor=magenta,      
    urlcolor=cyan,
}
\newcounter{teststepctr}


\definecolor{gray}{cmyk}{0.05,0.05,0.05,0.05}

\urlstyle{same}

%\setlength{\parindent}{4cm}

\pagestyle{fancy}
\renewcommand{\headrulewidth}{0pt}
\fancyhead[CE,CO,LE,LO,RE,RO]{} %% clear out all headers
\fancyhead[C]{%
           \begin{flushleft}
          \includegraphics[height=2.5cm,width=2.5cm]{Starfleet_Corps_of_Engineers_Logo.png}
          \end{flushleft}
          \bgroup
          \def\arraystretch{1.5}% 
          \begin{tabular}{V{4}l|l|l|lV{4}}
            \hlineB{4}
            \multicolumn{4}{V{4}cV{4}}{
                \textbf{Requirements and Test Specification - Self-Sealing Stem Bolt }
            } \\ 
            \hline
            {\footnotesize Type: Requirement \& Test Specification} & 
            {\footnotesize No.: 176-3049211-02} & 
            {\footnotesize Version: 1} & 
            {\footnotesize Page: \thepage of \pageref{LastPage}} \\ 
            \hlineB{4}
          \end{tabular}%
          \egroup
}


\begin{document}

\section{Purpose}
This document defines the hardware and software requirements as well as test procedures required for verification of a self-sealing stem bolt. As we know nothing could be built without bolts. They are a basic component of reverse-ratcheting routing planers.
\section{References}

\subsection{Internal References}

\begin{itemize}
	\item  Self-sealing stem bolt on memory-alpha.fandom.com (\href{https://memory-alpha.fandom.com/wiki/Self-sealing_stem_bolt}{link})
	\item  Self-sealing stem bolt on memory-beta.fandom.com (\href{https://memory-beta.fandom.com/wiki/Self-sealing_stembolt}{link})
\end{itemize}

\section{System Requirements}



\subsection{User Interface}

\begin{theorem}
	All system shall have an interface that allows users (operator, maintainer, engineer) to access the system directly.
\end{theorem}

\noindent\textbf{\large{Language and Labels}}

\begin{subtheorems}
	\begin{theorem}
		All labeling shall be in English or Andorian.
	\end{theorem}

	\begin{theorem}
		The part number, revision and serial number should be visible on the outside of the self-sealing stem bolt.
	\end{theorem}

	\noindent\textbf{\large{User Controller Parameters}}

	\begin{theorem}
		The user shall be able to select the magnetic flux capacitance flow through the self-sealing stem bolt by increments of 5 MEV and a range between 0 MeV and 1 TeV.
	\end{theorem}

	\begin{theorem}
		The user shall be able to read the current magnetic flux capacitance flow through the self-sealing stem bolt under the following light conditions: full sunlight, dusk, dawn, shadow, indoors, luminescent light, complete darkness.
	\end{theorem}

	\begin{theorem}
		The self-sealing stem bolt shall have a debug interface that allows authorized users access to
		low level functionality and operational logs.
	\end{theorem}


	\noindent\textbf{\large{Error Messages}}

	\begin{theorem}
		The self-sealing stem bolt shall have an audio-visual malfunction indicator.
	\end{theorem}
	\begin{theorem}
		The self-sealing stem bolt shall output detailed error messages via the debug interface in the case of an error.
	\end{theorem}


\end{subtheorems}

\subsection{Modes of Operation}


\begin{theorem}
	The self-sealing stem bolt shall have an operational mode and a maintenance mode.
\end{theorem}

\begin{subtheorems}

	\noindent\textbf{\large{Installation \& Startup}}


	\noindent\textbf{\large{Operational Mode}}


\end{subtheorems}




\section{System Verification}
This document describes the system verification for self-sealing stem bolt. Unfortunately self-sealing stem bolts are mysterious devices of unknown use and origin. The field of application of the self-sealing stem bolts is shrouded in mystery.

\subsection{Objective}

Even though the self-sealing stem bolt is a mysterious device, this document outlines a procedure to determine how one can be tested.

\subsection{Required Equipment}

\begin{itemize}
	\item standard issue Tricorder
	\item 20T NMR spectrometer
	\item quantum combobulator
\end{itemize}



\subsection{Setup and Configuration}

No special setup and configuration is required beyond ensuring calibration of Tricorder, spectrometer and combobulator.

\subsection{Requirements Tested}


\bgroup
\def\arraystretch{1.5}% 
\begin{tabularx}{\textwidth}{|X|}
	\hline
	\textbf{Requirements from Document Number:} \\ Reference the document where the requirements are coming from here \\ \hline
	S3.1.1, S3.2.2, S3.2.3, S3.2.4              \\ \hline
\end{tabularx}
\egroup

\subsubsection{Procedure and Test Worksheet}

\bgroup
\def\arraystretch{1.5}% 
\begin{tabularx}{\textwidth}{|>{\hsize=.8\hsize\linewidth=\hsize}X|>{\hsize=1.2\hsize\linewidth=\hsize}X|}
	\hline
	\textbf{Stardate Test Performed}           & \\ \hline
	\textbf{Starbase Test Performed}           & \\ \hline
	\textbf{Test System Serial \#}             & \\ \hline
	\textbf{Test Computer Serial \#}           & \\ \hline
	\textbf{Functional self-sealing stem bolt} & \\ \hline
	\textbf{Defective self-sealing stem bolt}  & \\ \hline
\end{tabularx}
\egroup




\bgroup

\small
\def\arraystretch{1.5}% 
\begin{xltabular}{\textwidth}{
	|>{\hsize=0.25\hsize\linewidth=\hsize}X
	|>{\hsize=1.85\hsize\linewidth=\hsize}X
	|>{\hsize=1.85\hsize\linewidth=\hsize}X
	|>{\hsize=0.25\hsize\linewidth=\hsize}X
	|>{\hsize=0.8\hsize\linewidth=\hsize}X
	|
	}
	\hline
	\rowcolor{gray}


	Step  &
	Action  &
	Verification  &
	P/F  &
	Result, Notes   \\ \hline \endhead
	\addtocounter{teststepctr}{1} \theteststepctr \ &
	Ensure indicator lights are working by pressing and holding down the on button for 10 seconds.&
	All indicator lights (red, gree, blue) flash first one after another twice and then all together 3 times.
	&

	&

	\\ \hline

	\addtocounter{teststepctr}{1} \theteststepctr \ &
	Ensure correct composition of self-sealing stem bolt by inserting it into a spectrometer &
	The stem bolt should consist of 80\% mixed duranium, aluminum, and steel alloys, 11\% electrically modulated ceramic, and 9\% thermally stabilized plastic.
	&

	&

	\\ \hline

	\addtocounter{teststepctr}{1} \theteststepctr \ &
	Insert self-stealing stem bolt in quantum combobulator and verify left and right-handed quantum flux capacitance &
	The measured flux capacitance on the left hand should not exceed an output conductance of 17.8 $\pm$ 2 millimohs.
	The measured flux capacitance on the right hand should not exceed an output conductance of 3.2 $\pm$ 1 nanomohs.
	&

	&

	\\ \hline

	\addtocounter{teststepctr}{1} \theteststepctr \ &
	Verify the positive and negative ion flow in upper and lower unilateral phase detractor alignment assemblies when applying a 750V with a phase offset of 5403 Furmans.  &
	The negative ion flow in the upper unilateral phase detractor alignment assembly does not exceed 56 Mol/Angstrom and 23 Mol/Anstrom in the lower alignment assembly.
	The measured flux capacitance on the right hand should not exceed an output conductance of 3.2 $\pm$ 1 nanomohs.
	&

	&

	\\ \hline

	\addtocounter{teststepctr}{1} \theteststepctr \ &
	Verify the positive and negative ion flow in upper and lower unilateral phase detractor alignment assemblies when applying a 750V with a phase offset of 14664 Furmans.  &
	The negative ion flow in the upper unilateral phase detractor alignment assembly does not exceed 104 Mol/Angstrom and 72 Mol/Anstrom in the lower alignment assembly.
	The measured flux capacitance on the right hand should not exceed an output conductance of 6.32 $\pm$ 2 nanomohs.
	&

	&

	\\ \hline

\end{xltabular}
\egroup


\noindent\bgroup
\def\arraystretch{1.5}% 
\begin{tabularx}{\textwidth}{|X|}
	\hline
	\textbf{Summary Report Results:} \\  \hline
	Additional Comments:             \\ \rule{0pt}{12ex}     \\ \hline


	\bgroup
	\def\arraystretch{1.1}% 
	\begin{tabular}{@{}p{0.98in}p{3.2in}p{0.3in}p{1.4in}@{}}
		\rule{0pt}{3ex}                                                  \\
		Completed by: & \hrulefill                  & Date: & \hrulefill \\
		              & Lieutenant junior grade Nog &       &            \\
		              & Junior Engineer             &       &            \\
	\end{tabular}
	\egroup
	\\ \hline


	\bgroup
	\def\arraystretch{1.1}% 
	\begin{tabular}{@{}p{0.98in}p{1.45in}p{1.45in}p{1.6in}@{}}
		\rule{0pt}{.5ex}                                 \\ Result: &
		Pass $\square$ &
		Fail $\square$ &
		Pass with limitations $\square$ \rule{0pt}{.5ex} \\
	\end{tabular}
	\egroup
	\\ \hline



	\bgroup
	\def\arraystretch{1.1}% 
	\begin{tabular}{@{}p{0.98in}p{3.2in}p{0.3in}p{1.4in}@{}}
		\rule{0pt}{3ex}                                                          \\
		Approved By: & \hrulefill                           & Date: & \hrulefill \\
		             & Lieutenant commander Geordi La Forge &       &            \\
		             & Chief Engineer                       &       &            \\
	\end{tabular}
	\egroup
	\\ \hline
\end{tabularx}
\egroup


\section{Document Revision History}



\bgroup

\small
\def\arraystretch{1.5}% 
\begin{xltabular}{\textwidth}{
	|>{\hsize=0.15\hsize\linewidth=\hsize}X
	|>{\hsize=1.7\hsize\linewidth=\hsize}X
	|>{\hsize=0.15\hsize\linewidth=\hsize}X
	|
	}
	\hline
	\rowcolor{gray}


	Ver  &
	Change Description  &
	Author  \\ \hline \endhead

	01 & Release version 1 of self-sealing stem bolt & Nog \\ \hline

\end{xltabular}
\egroup

\end{document}
